\BeleskeID{08-s006}
% element: 08-s006:e04
% element: 08-s006:notes
Povezivanje između kola i prema izvoru izvodi se žicama, metalnim linijama na štampanoj ploči ili metalnim linijama unutar integrisanog kola. Takve veze nisu idealne: imaju otpornost i parazitne efekte i u sebi i prema okolini, uključujući kapacitivnosti i induktivnosti.

Napon $b_1(t)$ meri se prema lokalnoj masi prvog kola, a $a_2(t)$ prema lokalnoj masi drugog. Struja povratnog voda stvara pad napona preko $Z_p$, pa isti signalni vod ne garantuje isti ulazni napon prema obe lokalne mase. Smanjenje impedanse povratnog puta zato smanjuje degradaciju. Put treba da bude što kraći, sa što manje rupa i uskih grla, uz što manju otpornost i parazitne efekte.
